\chapter{Jumlah Peubah Acak dan Hukum Bilangan Besar}\label{ch:g12:sums}

Mengapa kasino selalu menang pada akhirnya, dan mengapa jajak pendapat berhasil? Sebab
rata-rata banyak besaran acak yang \omterm{def:g12:sums:indep}{saling bebas} semakin sedikit berayun. Bab
ini membuktikannya: kelinearan harapan, keaditifan \omterm{def:g12:randvar:exp}{ragam} bagi peubah
yang \omterm{def:g12:sums:indep}{saling bebas}, ketaksamaan Bienaymé--Chebyshev, dan hukum bilangan
besar.

\section{Jumlah peubah acak}

Diberikan dua \omterm{def:g12:randvar:rv}{peubah acak} $X, Y$ pada \omterm{def:g11:prob:model}{ruang sampel} berhingga yang sama,
$\Omega$, jumlah $X + Y$ adalah \omterm{def:g12:randvar:rv}{peubah acak}
$\omega \mapsto X(\omega) + Y(\omega)$.

\begin{theorem}[Kelinearan harapan]\label{thm:g12:sums:linearity}
\index{harapan!kelinearan}
Untuk semua \omterm{def:g12:randvar:rv}{peubah acak} $X, Y$ pada $\Omega$ dan $a, b \in \R$:
\[
\E(X + Y) = \E(X) + \E(Y), \qquad \E(aX + b) = a\E(X) + b .
\]
\emph{Tak ada anggapan \omterm{def:g12:condprob:indep}{kebebasan} yang diperlukan.}
\end{theorem}

\begin{proof}
Tulislah harapannya sebagai jumlah atas hasilnya: karena
$\P(X = x) = \sum_{\omega : X(\omega) = x} \P(\{\omega\})$, mengelompokkan sukunya
memberi $\E(X) = \sum_{\omega \in \Omega} \P(\{\omega\})\,X(\omega)$. Maka
\[
\E(X + Y) = \sum_{\omega} \P(\{\omega\})\bigl(X(\omega) + Y(\omega)\bigr)
= \sum_{\omega} \P(\{\omega\})X(\omega)
+ \sum_{\omega} \P(\{\omega\})Y(\omega) = \E(X) + \E(Y). \qedhere
\]
\end{proof}

\begin{definition}[Peubah acak yang saling bebas]\label{def:g12:sums:indep}
Peubah $X$ dan $Y$ disebut \emph{saling bebas}\index{peubah acak!saling bebas} jika
untuk semua nilai $x, y$:
\[
\P(X = x \text{ dan } Y = y) = \P(X = x)\,\P(Y = y).
\]
Beberapa peubah $X_1, \dots, X_n$ disebut saling bebas jika kaidah hasil kali ini
berlaku bagi setiap pilihan nilai dari setiap subkeluarganya.
\end{definition}

\begin{proposition}\label{prop:g12:sums:prodexp}
Jika $X$ dan $Y$ \omterm{def:g12:sums:indep}{saling bebas}, maka $\E(XY) = \E(X)\,\E(Y)$.
\end{proposition}

\begin{proof}
\[
\begin{aligned}
\E(XY) &= \sum_{x, y} xy\;\P(X = x \text{ dan } Y = y)
= \sum_{x, y} xy\,\P(X=x)\P(Y=y) \\
&= \Bigl(\sum_x x\P(X=x)\Bigr)\Bigl(\sum_y y\P(Y=y)\Bigr). \qedhere
\end{aligned}
\]
\end{proof}

\begin{theorem}[Ragam sebuah jumlah]\label{thm:g12:sums:variance}
\index{ragam!sebuah jumlah}
Jika $X$ dan $Y$ \emph{\omterm{def:g12:sums:indep}{saling bebas}}, maka
\[
\V(X + Y) = \V(X) + \V(Y).
\]
Secara lebih umum, untuk $X_1, \dots, X_n$ yang \omterm{def:g12:sums:indep}{saling bebas}:
$\V(X_1 + \dots + X_n) = \V(X_1) + \dots + \V(X_n)$.
\end{theorem}

\begin{proof}
Dengan memakai König--Huygens (\cref{prop:g12:randvar:konig}) dan kelinearannya:
\begin{align*}
\V(X + Y) &= \E\bigl((X+Y)^2\bigr) - \bigl(\E X + \E Y\bigr)^2\\
&= \E(X^2) + 2\E(XY) + \E(Y^2)
- \E(X)^2 - 2\E(X)\E(Y) - \E(Y)^2\\
&= \V(X) + \V(Y) + 2\bigl(\E(XY) - \E(X)\E(Y)\bigr),
\end{align*}
dan kurung terakhirnya lenyap bagi peubah yang \omterm{def:g12:sums:indep}{saling bebas}
(\cref{prop:g12:sums:prodexp}). Kasus umumnya menyusul lewat induksi.
\end{proof}

\begin{example}[Binomial ditinjau ulang]\label{ex:g12:sums:binomial}
Peubah binomial $X \sim \mathcal B(n, p)$ adalah jumlah
$X = X_1 + \dots + X_n$ berisi $n$ peubah Bernoulli yang \omterm{def:g12:sums:indep}{saling bebas}. Karena itu,
secara berstruktur:
\[
\E(X) = np, \qquad \V(X) = np(1-p),
\]
sehingga \cref{thm:g12:randvar:binomial} diperoleh kembali tanpa perhitungan apa pun.
\end{example}

\begin{proposition}[Rata-rata contoh]\label{prop:g12:sums:mean}
Misalkan $X_1, \dots, X_n$ \omterm{def:g12:sums:indep}{saling bebas} dan berdistribusi sama dengan $X$
(harapannya $\mu$, \omterm{def:g12:randvar:exp}{ragamnya} $\sigma^2$), lalu misalkan
$M_n = \frac{X_1 + \dots + X_n}{n}$ adalah \emph{rata-rata
contohnya}\index{rata-rata contoh}. Maka
\[
\E(M_n) = \mu, \qquad \V(M_n) = \frac{\sigma^2}{n}, \qquad
\sigma(M_n) = \frac{\sigma}{\sqrt n}.
\]
\end{proposition}

\begin{proof}
Kelinearannya memberi $\E(M_n) = \frac{n\mu}{n} = \mu$; \omterm{def:g12:condprob:indep}{kebebasannya} memberi
$\V(X_1 + \dots + X_n) = n\sigma^2$, dan pembagian dengan $n$ menskalakan
\omterm{def:g12:randvar:exp}{ragamnya} dengan $\frac{1}{n^2}$ (\cref{prop:g12:randvar:affine}).
\end{proof}

Bilangan $\frac{\sigma}{\sqrt n}$ itulah \emph{hukum akar kuadrat} yang mendasar: untuk
menyeparuhkan ayunan sebuah rata-rata, lipatempatkanlah ukuran contohnya.

\begin{omfigure}
\begin{tikzpicture}
\begin{axis}[omaxis,
  width=9.6cm, height=5.6cm,
  xlabel={$n$}, ylabel={$\sigma(M_n)/\sigma$},
  xmin=0, xmax=105, ymin=0, ymax=1.15,
  xtick={4,16,64}, ytick={1,0.5,0.25},
  yticklabels={$1$,$\tfrac12$,$\tfrac14$}]
  \addplot[omDef, thick, domain=1:100, samples=80] {1/sqrt(x)};
  \draw[black, dashed, thin] (axis cs:0,0.5) -- (axis cs:4,0.5)
    -- (axis cs:4,0);
  \draw[black, dashed, thin] (axis cs:0,0.25) -- (axis cs:16,0.25)
    -- (axis cs:16,0);
  \draw[black, dashed, thin] (axis cs:64,0.125) -- (axis cs:64,0);
  \fill (axis cs:1,1) circle (1.5pt);
  \fill (axis cs:4,0.5) circle (1.5pt);
  \fill (axis cs:16,0.25) circle (1.5pt);
  \fill (axis cs:64,0.125) circle (1.5pt);
\end{axis}
\end{tikzpicture}

{\small Hukum akar kuadrat: tiap kali ukuran contohnya dilipatempatkan, \omterm{def:g11:stat:variance}{simpangan
baku} rata-ratanya hanya terpotong separuh.}
\end{omfigure}

\section{Ketaksamaan pemusatan}

\begin{theorem}[Ketaksamaan Markov]\label{thm:g12:sums:markov}
\index{ketaksamaan Markov}
Jika $X \geq 0$ dan $a > 0$:
\[
\P(X \geq a) \leq \frac{\E(X)}{a}.
\]
\end{theorem}

\begin{proof}
Pada $\E(X) = \sum_i p_i x_i$ (semua sukunya tak negatif), pertahankan hanya suku
dengan $x_i \geq a$: masing-masing sekurang-kurangnya $a\,p_i$, sehingga
$\E(X) \geq a \sum_{x_i \geq a} p_i = a\,\P(X \geq a)$.
\end{proof}

\begin{theorem}[Ketaksamaan Bienaymé--Chebyshev]\label{thm:g12:sums:chebyshev}
\index{ketaksamaan Bienaymé--Chebyshev}
Untuk sebarang \omterm{def:g12:randvar:rv}{peubah acak} $X$ dan sebarang $\varepsilon > 0$:
\[
\P\bigl(\abs{X - \E(X)} \geq \varepsilon\bigr)
\leq \frac{\V(X)}{\varepsilon^2}.
\]
\end{theorem}

\begin{proof}
Terapkanlah ketaksamaan Markov pada peubah tak negatif
$Y = (X - \E(X))^2$ dengan $a = \varepsilon^2$:
\[
\P\bigl(\abs{X - \E(X)} \geq \varepsilon\bigr)
= \P(Y \geq \varepsilon^2)
\leq \frac{\E(Y)}{\varepsilon^2} = \frac{\V(X)}{\varepsilon^2}. \qedhere
\]
\end{proof}

\begin{omfigure}
\begin{tikzpicture}
\begin{axis}[omaxis,
  width=9.6cm, height=5cm,
  ylabel={}, xmin=-3.4, xmax=3.4, ymin=0, ymax=1.25,
  xtick={-1,0,1},
  xticklabels={$\E(X)-\varepsilon$,$\E(X)$,$\E(X)+\varepsilon$},
  ytick=\empty]
  \addplot[name path=bell, omDef, thick, domain=-3.2:3.2]
    {exp(-x^2/2)};
  \path[name path=xaxis] (-3.2,0) -- (3.2,0);
  \addplot[omDef!20] fill between[of=bell and xaxis,
    soft clip={domain=-1:1}];
  \addplot[omThm!20] fill between[of=bell and xaxis,
    soft clip={domain=-3.2:-1}];
  \addplot[omThm!20] fill between[of=bell and xaxis,
    soft clip={domain=1:3.2}];
  \draw[black, dashed, thin] (axis cs:-1,0) -- (axis cs:-1,0.75);
  \draw[black, dashed, thin] (axis cs:1,0) -- (axis cs:1,0.75);
\end{axis}
\end{tikzpicture}

{\small Pemusatannya: Bienaymé--Chebyshev membatasi peluang bahwa
$X$ jatuh pada ekornya (merah), pada jarak sekurang-kurangnya $\varepsilon$ dari
harapannya, oleh $\V(X)/\varepsilon^2$.}
\end{omfigure}

\begin{theorem}[Hukum bilangan besar]\label{thm:g12:sums:lln}
\index{hukum bilangan besar}
Misalkan $X_1, \dots, X_n$ \omterm{def:g12:sums:indep}{saling bebas} dan berdistribusi sama
(harapannya $\mu$, \omterm{def:g12:randvar:exp}{ragamnya} $\sigma^2$), dan $M_n$ \omterm{prop:g12:sums:mean}{rata-rata contohnya}. Untuk
setiap $\varepsilon > 0$:
\[
\P\bigl(\abs{M_n - \mu} \geq \varepsilon\bigr)
\leq \frac{\sigma^2}{n\,\varepsilon^2}
\xrightarrow[n \to +\infty]{} 0 .
\]
\omterm{prop:g12:sums:mean}{Rata-rata contohnya} memusat di sekeliling harapannya.
\end{theorem}

\begin{proof}
Bienaymé--Chebyshev yang diterapkan pada $M_n$, yang harapannya $\mu$ dan
\omterm{def:g12:randvar:exp}{ragamnya} $\frac{\sigma^2}{n}$ (\cref{prop:g12:sums:mean}).
\end{proof}

\begin{remark}
Teorema ini adalah jembatan antara teori peluang dan statistika:
\emph{frekuensi} sebuah \omterm{def:g11:prob:model}{kejadian} pada banyak ulangan yang \omterm{def:g12:sums:indep}{saling bebas} mendekati
\emph{peluangnya} (ambillah $X_i$ sebagai penunjuk \omterm{def:g11:prob:model}{kejadiannya}, sehingga
$\mu = p$). Teorema itu membenarkan penaksiran peluang lewat simulasi
(\emph{metode Monte Carlo}) dan penaksiran proporsi populasi lewat jajak pendapat ---
secara kuantitatif: lihat \cref{ch:g12:contdist}.
\end{remark}

\begin{method}[Memakai Bienaymé--Chebyshev]\label{met:g12:sums:chebyshev}
Untuk menjamin $\P(\abs{M_n - \mu} \geq \varepsilon) \leq \alpha$,
cukuplah $n \geq \frac{\sigma^2}{\alpha\,\varepsilon^2}$. Batasnya
kasar (ayunan yang sesungguhnya biasanya jauh lebih kecil) tetapi berlaku umum
sepenuhnya: batas itu tak menuntut apa pun tentang \omterm{def:g12:randvar:rv}{distribusinya} selain \omterm{def:g12:randvar:exp}{ragam} yang berhingga.
\end{method}

\section{Latihan}

\begin{exercise}[$\star$]\label{exo:g12:sums:1}
Dua dadu setimbang dilempar; misalkan $S$ jumlahnya. Dengan memakai kelinearannya (bukan
\omterm{def:g12:randvar:rv}{distribusi} $S$!), hitunglah $\E(S)$; lalu hitung $\V(S)$ dengan memakai
\omterm{def:g12:condprob:indep}{kebebasannya}, jika diketahui bahwa satu dadu setimbang beragam $\frac{35}{12}$.
\end{exercise}

\begin{exercise}[$\star$]\label{exo:g12:sums:2}
Misalkan $X \sim \mathcal B(100,\ 0.5)$. Batasilah
$\P(X \geq 75)$ dengan ketaksamaan Markov, lalu
$\P(\abs{X - 50} \geq 25)$ dengan Bienaymé--Chebyshev. Bandingkan.
\end{exercise}

\begin{exercise}[$\star$]\label{exo:g12:sums:3}
Sekeping uang logam yang setimbang dilempar $n$ kali dan $F_n$ menyatakan frekuensi gambarnya.
Seberapa besar $n$ harus dibuat agar, menurut batas Bienaymé--Chebyshev,
$\P\bigl(\abs{F_n - 0.5} \geq 0.05\bigr) \leq 0.05$?
\end{exercise}

\begin{exercise}[$\star\star$]\label{exo:g12:sums:4}
Seorang pemodal membelah modalnya sama rata di antara $n$ aset yang \omterm{def:g12:sums:indep}{saling bebas},
masing-masing berharapan imbal hasil $\mu = 5\%$ dan bersimpangan baku
$\sigma = 20\%$. Hitunglah harapan dan \omterm{def:g11:stat:variance}{simpangan baku}
imbal hasil portofolionya $M_n$, lalu banyaknya aset yang diperlukan untuk membawa
\omterm{def:g11:stat:variance}{simpangan bakunya} di bawah $4\%$. Asas keuangan apa yang
digambarkan hal ini?
\end{exercise}

\begin{exercise}[$\star\star$]\label{exo:g12:sums:5}
Misalkan $X$ dan $Y$ \omterm{def:g12:sums:indep}{saling bebas}, dan keduanya seragam pada $\{1, 2, 3\}$.
\begin{enumerate}
  \item Berikan \omterm{def:g12:randvar:rv}{distribusi} $S = X + Y$ lalu hitunglah $\E(S)$, $\V(S)$
        langsung darinya.
  \item Perolehlah kembali kedua nilainya lewat kelinearan dan keaditifan \omterm{def:g12:randvar:exp}{ragamnya}.
\end{enumerate}
\end{exercise}

\begin{exercise}[$\star\star$]\label{exo:g12:sums:6}
Tunjukkan bahwa keaditifan \omterm{def:g12:randvar:exp}{ragamnya} dapat gagal tanpa \omterm{def:g12:condprob:indep}{kebebasan}:
hitunglah $\V(X + Y)$ untuk $Y = X$, lalu bandingkan dengan $\V(X) + \V(Y)$. Untuk
peubah $X$ yang mana berlaku $\V(2X) = 2\V(X)$?
\end{exercise}

\begin{exercise}[$\star\star$]\label{exo:g12:sums:7}
Sebuah dadu dicurigai berat sebelah. Dadu itu dilempar $1200$ kali dan menunjukkan
angka enam $260$ kali ($f = 0.2167$ alih-alih $\frac16 \approx 0.1667$). Di bawah
hipotesis bahwa dadunya setimbang, batasilah
$\P\bigl(\abs{F_n - \frac16} \geq 0.05\bigr)$ dengan Bienaymé--Chebyshev, lalu
bahaslah apakah hipotesis kesetimbangannya masuk akal.
\end{exercise}

\begin{exercise}[$\star\star\star$]\label{exo:g12:sums:8}
\emph{(Ketaksamaan yang lebih baik bagi uang logamnya.)} Misalkan
$X \sim \mathcal B(n,\ p)$ dan $F_n = \frac Xn$.
\begin{enumerate}
  \item Tunjukkan bahwa $p(1-p) \leq \frac14$ untuk $p \in \intcc{0}{1}$.
  \item Simpulkan batas yang bebas \omterm{def:g12:randvar:rv}{distribusi},
        $\P\bigl(\abs{F_n - p} \geq \varepsilon\bigr) \leq
        \frac{1}{4n\varepsilon^2}$.
  \item Berapa banyak orang yang harus dijajaki agar frekuensi yang teramati berada
        dalam jarak $3$ angka dari proporsi yang sebenarnya dengan peluang sekurang-kurangnya
        $95\%$, menurut batas ini? (Jajak pendapat yang sungguhan memakai taksiran yang lebih tajam, tetapi
        orde besarnya sudah tepat.)
\end{enumerate}
\end{exercise}

\section{Soal: Bandar selalu menang}

\begin{problem}[{Soal akhir pekan --- hukum bilangan besar
menjelaskan kasino, jajak pendapat dan asuransi, lalu meruntuhkan
sesat pikir penjudi di sepanjang jalannya}]
\label{pb:g12:sums:1}
Seorang pemain rolet sesudah $100$ putaran, hampir sesering tidak,
sedang \emph{unggul}. Kasino yang menjalankan sejuta putaran unggul dengan
kepastian yang takkan dipersoalkan pengadilan mana pun. Permainan yang sama, keunggulan
mungil yang sama --- \omterm{def:g11:seq:arithmetic}{bedanya} adalah $\sqrt n$, dan itulah pokok bab
ini (\cref{prop:g12:sums:mean},
\cref{thm:g12:sums:chebyshev}, \cref{thm:g12:sums:lln}). Soal
ini menjalankan pembukuan kasinonya, mengukur jajak pendapat pemilu, menetapkan harga
penganekaragaman --- lalu membongkar sesat pikir termahal dalam
sejarah perjudian.

\medskip
\noindent\textbf{Bagian I --- Kelancaran.}
\begin{enumerate}
  \item Dua dadu yang \omterm{def:g12:sums:indep}{saling bebas}: hitunglah $V(X + Y)$
        (\cref{thm:g12:sums:variance}).
  \item Lemparlah $100$ dadu yang \omterm{def:g12:sums:indep}{saling bebas} lalu rata-ratakan hasilnya:
        berikan harapan dan \omterm{def:g11:stat:variance}{simpangan baku}
        \omterm{prop:g12:sums:mean}{rata-rata contohnya} (\cref{prop:g12:sums:mean}; untuk satu dadu,
        $\sigma = \sqrt{35/12} \approx 1.71$).
  \item Sebuah \omterm{def:g12:randvar:rv}{peubah acak} tak negatif berharapan $2$.
        Apa yang dikatakan ketaksamaan Markov
        (\cref{thm:g12:sums:markov}) tentang
        $\P(X \geq 10)$?
  \item Batasilah $\P\left(\abs{\bar X_{100} - 3.5} \geq
        0.5\right)$ bagi rata-rata dadu pada pertanyaan 2 dengan
        Bienaymé--Chebyshev.
  \item Harapan selalu menjumlah, sedangkan \omterm{def:g12:randvar:exp}{ragam} hanya menjumlah di bawah
        \omterm{def:g12:condprob:indep}{kebebasan}: tunjukkanlah $X, Y$ yang bergantungan (petunjuk:
        $Y = -X$) yang membuat
        $V(X + Y) \neq V(X) + V(Y)$.
\end{enumerate}

\medskip
\noindent\textbf{Bagian II --- Pembukuan sang bandar.}
Rolet Eropa, bertaruh $1$ euro pada merah: menang $+1$ dengan
peluang $\frac{18}{37}$, kalah $-1$ dengan peluang
$\frac{19}{37}$.
\begin{enumerate}[resume]
  \item Hitunglah harapan dan \omterm{def:g11:stat:variance}{simpangan baku} keuntungan satu
        taruhan.
  \item Seorang penjudi memasang $100$ taruhan yang \omterm{def:g12:sums:indep}{saling bebas}; misalkan $G$
        keuntungan totalnya. Hitunglah $\E(G)$ dan $\sigma(G)$, lalu
        nisbah $\frac{\abs{\E(G)}}{\sigma(G)}$. Apa arti
        hanyutan seperempat \omterm{def:g11:stat:variance}{simpangan baku} bagi
        peluang penjudinya untuk unggul malam itu?
  \item Kasinonya melihat $1\,000\,000$ taruhan. Hitunglah
        harapan dan \omterm{def:g11:stat:variance}{simpangan baku} pendapatan totalnya,
        lalu nisbah yang sama. Tafsirkan bilangan $27$ itu.
  \item Sahkanlah dengan Chebyshev: batasilah peluang bahwa
        kasinonya \emph{rugi} uang sepanjang sejuta taruhan itu.
  \item Sistem bertaruh: melipatduakan sesudah kalah, berhenti ketika
        unggul\,\dots\ Dengan memakai kelinearan harapan
        (\cref{thm:g12:sums:linearity}) atas \omterm{def:g12:seq:sequence}{barisan} taruhannya (yang
        mungkin acak), jelaskan mengapa \emph{setiap}
        siasat pada permainan berkeunggulan negatif berharapan keuntungan
        negatif --- apa yang akan dilanggar oleh sistem yang menang?
  \item Nyatakan apa yang dijanjikan hukum bilangan besar
        (\cref{thm:g12:sums:lln}) tentang frekuensi
        merahnya --- dan apa yang \emph{tidak} dijanjikannya tentang
        putaran berikutnya sesudah sepuluh merah berturut-turut. Sebutkan
        sesat pikirnya.
  \item Butir yang paling halus: hukumnya bekerja lewat \emph{pengenceran},
        bukan lewat penggantian. Jika merahnya berjalan $100$ di atas
        harapannya sesudah suatu malam, kelebihan itu tak pernah
        ``dibayar kembali'' --- hitunglah apa yang justru terjadi pada
        \emph{pecahan} $\frac{100}{n}$ ketika $n = 10^6$, lalu
        tulislah ulang kekeliruan sesat pikir penjudi itu dalam satu kalimat.
\end{enumerate}

\medskip
\noindent\textbf{Bagian III --- Jajak pendapat.}
\begin{enumerate}[resume]
  \item Dari \cref{exo:g12:sums:8}: berapa banyak orang yang harus
        dijajaki agar frekuensi yang teramati jatuh dalam jarak $3$
        angka dari kebenarannya dengan peluang sekurang-kurangnya
        $95\,\%$, menurut batas Chebyshev yang bebas \omterm{def:g12:randvar:rv}{distribusi} itu?
  \item Lembaga jajak pendapat yang sungguhan memakai kira-kira $1\,100$ orang untuk
        marjin $\pm 3$ angka pada $95\,\%$: rumusnya adalah
        $n \approx \frac{1.96^2 \times p(1-p)}
        {\varepsilon^2}$ dengan $p(1-p) \leq \frac14$. Hitunglah
        nilainya, lalu jelaskan jurangnya dengan pertanyaan 13 (apa
        yang tidak diketahui Chebyshev tentang bentuk
        ayunannya?).
  \item Untuk menyeparuhkan marjin galatnya, bagaimana contohnya harus
        bertambah? Dari $\pm 3$ angka pada $1\,100$, contoh berapa
        yang memberi $\pm 1.5$ angka?
  \item Klasik yang berlawanan dengan naluri: ukuran contoh yang diperlukan
        tak pernah memakai ukuran populasinya --- $1\,100$ orang
        cukup bagi sebuah kota maupun sebuah benua. Tunjuklah tempat
        di dalam modelnya yang di sana ukuran populasinya tak hadir, lalu
        berikan perumpamaan dapur yang dipakai para penjajak pendapat.
\end{enumerate}

\medskip
\noindent\textbf{Bagian IV --- Penganekaragaman.}
\begin{enumerate}[resume]
  \item Penanggung asuransi pada \cref{pb:g12:randvar:1} memperoleh
        $30n$ euro sebagai harapan atas $n$ polis dengan
        \omterm{def:g11:stat:variance}{simpangan baku} klaim kira-kira $315\sqrt n$. Untuk
        $n$ yang mana harapan labanya akhirnya melampaui satu
        \omterm{def:g11:stat:variance}{simpangan baku} klaimnya? Apa yang sedang dikerjakan
        $\sqrt n$ bagi penanggung asuransinya?
  \item Imbal hasil tahunan satu saham berayun dengan
        $\sigma = 20\,\%$. Belahlah uangnya sama rata atas $25$
        saham semacam itu yang \emph{\omterm{def:g12:sums:indep}{saling bebas}}: hitunglah
        $\sigma$ portofolionya. Keuangan menyebut penganekaragaman
        satu-satunya makan siang gratis --- apa tepatnya makan siang itu?
  \item Sekarang biarkan $25$ sahamnya berkorelasi sempurna (semuanya
        bergerak bersama): berapa $\sigma$ portofolionya?
        Bandingkan kedua ujungnya lalu nyatakan anggapan mana
        yang diam-diam disewa setiap alasan penganekaragaman ---
        dan apa yang terjadi ketika anggapan itu gagal di seluruh sistem pada tahun 2008.
  \item Penutup --- simfoni $\sqrt n$: jumlahnya berayun seperti
        $\sqrt n$ sementara rata-ratanya mengendap seperti
        $\frac{1}{\sqrt n}$; mainkanlah temanya lewat keempat
        industri pada soal ini (kasino, jajak pendapat, asuransi,
        portofolio), sebutkan kedua penyalahguna yang tetap ada (sesat
        pikir penjudi dan \omterm{def:g12:condprob:indep}{kebebasan} yang palsu), lalu berikan
        penunjuk ke depan: \emph{bentuk} ayunannya
        --- yaitu loncengnya --- adalah bintang sinambung bab
        berikutnya, dan teorema lengkapnya memahkotai jilid
        universitas.

\end{enumerate}
\end{problem}
